\renewcommand{\chaptername}{Protocolo}
%==========================================
\chapter{}

% \begin{flushleft}
% 	Ing. Angel Eduardo Gómez Rodríguez\\
% 	Director, Dr. Dunstano del Puerto Flores \\
% 	Co-Director, Dr. \\
% 	Asesor, Dr. 
% \end{flushleft}

\section{Introducción}
Desde finales del sigolo XIX, Thomas A. Edison comenzo con el trabjo de electrificacion de luz electrica,
creando por primera vez, una red electrica de corriente continua (CC) en New York en 1882, la estacion Pearl Street 
generaba 30 kW a 110 Volts para un total de 59 clientes. Luego, Nikola desarrolo papers sobre las ventjas 
de la corriente alterna (CA), que describian motores de induccion y sincronos de dos fases. En 1891 se
puso en marcha la primera red electrica trifasica de CA en Alemnaia, transmitinendo energia a 179 kM a 12 kV
Para estados unidos en 1893 se puso la priemra red eelectrica de CA en California, uan red de 12 kM a 2.3 kV. 
En Mexico, la primera red electrica se deribo de una fabrica de tejidos en la ciudad de Leon, Gunajutato 
en 1879.

Los Sistemas Eléctricos de Potencia (SEP) modernos representan una de las infraestructuras más complejas 
y críticas de la sociedad contemporánea. A lo largo de las décadas, estas redes han experimentado una 
expansión significativa en escala e interconexión. La operación segura y fiable de estos vastos sistemas
 interconectados depende intrínsecamente de su capacidad para mantener la \textbf{estabilidad}~
 \cite{2}. Esta se define como la habilidad del sistema para retornar a un estado de 
 equilibrio operativo después de ser sujeto a una perturbación. La pérdida de estabilidad puede conducir 
 a oscilaciones de potencia no amortiguadas, resultando en fallos en cascada y apagones a gran escala.

Este desafío se ha intensificado notablemente. La creciente \textbf{complejidad} de las redes ya no 
se debe solo a su tamaño, sino a la integración masiva de fuentes de energía renovable basadas en 
convertidores, como la eólica y la solar, así como al uso de electrónica de potencia (HVDC, 
FACTS). Estos dispositivos introducen dinámicas fundamentlmente diferentes a la tradicional, 
caracterizadas por una baja inercia y controles rápidos. Esta transformación está impactando el 
amortiguamiento del sistema, haciendo al SEP más susceptible a oscilaciones electromecánicas de 
baja frecuencia (LFOs)~\cite{6} e incluso introduciendo nuevos fenómenos como las oscilaciones 
subsíncronas (SSO)~\cite{9}.

Ante esta creciente complejidad, la necesidad de un \textbf{monitoreo dinámico} en tiempo real es más 
crucial que nunca. Los sistemas tradicionales de supervisión (SCADA) son insuficientes para el análisis 
dinámico, debido a sus bajas tasas de muestreo y la falta de sincronización temporal entre mediciones 
distantes. La verdadera revolución tecnológica en este ámbito ha sido el desarrollo de las Unidades de 
Medición Fasorial (PMU) y su integración en Sistemas de Monitoreo de Área Amplia 
(WAMS)~\cite{5}. Las PMUs proporcionan mediciones de alta velocidad (e.g., 30-60 Hz) de 
fasores de voltaje y corriente, crucialmente sincronizadas mediante una marca de tiempo global (vía GPS),
 ofreciendo una análisis coherente del estado dinámico del sistema~\cite{3}.

La disponibilidad de flujos de datos de WAMS ha abierto dos vías principales para el análisis modal: 
el análisis de \textit{ringdown} tras una perturbación severa, y el análisis de \textbf{datos ambientales} 
(\textit{ambient data}) durante la operación normal~\cite{4, 8}. Mientras que las 
señales de \textit{ringdown} ofrecen una alta relación señal-ruido (SNR), son eventos infrecuentes. 
En contraste, los datos ambientales permiten un monitoreo continuo de la estabilidad. 
El desafío fundamental de los datos ambientales es que los modos de oscilación de interés están 
inherentemente ocultos bajo un nivel significativo de \textbf{ruido}  de medición, 
proveniente de las fluctuaciones aleatorias de las cargas y la generación~\cite{1}.

Esta realidad ha impulsado una intensa investigación en métodos de identificación modal basados en 
mediciones (\textit{data-driven}) capaces de operar en condiciones de bajo SNR. En la literatura se 
encuentran diversos enfoques, como el método de Prony~\cite{6, 9}, el Algoritmo 
de Realización del Ecosistema (ERA)~\cite{7, 6}, el método de Lápiz Matricial 
(Matrix Pencil)~\cite{6}, y el Ajuste Vectorial (Vector Fitting)~\cite{5}. 
Sin embargo, métodos clásicos como Prony son notoriamente sensibles al ruido, lo que degrada la precisión
 en la estimación de parámetros críticos como el amortiguamiento. Existe una \textbf{brecha} en la 
 búsqueda de un método que equilibre robustez al ruido, precisión en la estimación y la capacidad de 
 procesar de manera coherente la naturaleza \textbf{multi-señal} de los datos WAMS~\cite{10}.



 \section{Antecedentes} 



\section{Justificación}


 
\section{Objetivos}

\subsection{General}

\subsection{Específicos}
\begin{itemize}
	\item Analizar las técnicas de modelado y los modelos aplicables a las baterías de iones de litio que se ajusten a la pruebas de EIS. 
	\item Analizar las topologías de los convertidores de eletrónica de potencia que están presentes en los \emph{BESS} y puedan realizar la prueba de EIS.
	\item Simular y diseñar el sistema de control para realizar la prueba de EIS con el convertidor electrónico de potencia.
	  \item Aplicar los métodos de filtrado para validar el modelo de circuito eléctrico equivalente.
\end{itemize}

\section{Hipótesis}

\section{Metodología}


\begin{itemize}
    \item Revisión bibliográfica relacionada con:
\begin{itemize}
    \item Modelado de baterías de iones de litio: Se explorarán los modelos de circuito eléctrico equivalente y sus parámetros característicos, así como su relación con la impedancia electroquímica.
    \item Prueba de Espectroscopía de Impedancia Electroquímica (EIS): Se identificarán los métodos convencionales y las estrategias para aplicarlas mediante convertidores electrónicos de potencia.
    \item Topologías de convertidores electrónico de potencia: Se analizarán los convertidores más utilizados en los BESS que permitan la inyección de señales de prueba para la EIS.
\end{itemize}

\item Diseño y simulación del sistema de prueba.
\begin{itemize}
    \item Seleccionar la topología del convertidor electrónico de potencia para la implementación de la EIS.
    \item Desarrollar un modelo en software de simulación (MATLAB/Simulink) para evaluar el comportamiento del convertidor.
    \item Diseñar el sistema de control y medición encargado de modular el ciclo de trabajo del convertidor para la inyección de señales de prueba.
\end{itemize}

\item Implementación y experimentos.
\begin{itemize}
    \item Utilizar el convertidor electrónico de potencia para la realización de la prueba EIS a la batería de iones de litio modelo 18650.
    \item Ejecutar pruebas de EIS utilizando el convertidor electrónico de potencia como equipo de medición, inyectando perturbaciones controladas y registrando la respuesta de la batería.
\end{itemize}

\item Aplicación de métodos de filtrado y validación del modelo.
\begin{itemize}
    \item Aplicar técnicas de filtrado para procesar los datos obtenidos con el convertidor electrónico de potencia.
    \item Ajustar el modelo de circuito eléctrico equivalente de la batería en función de los datos experimentales obtenidos.
\end{itemize}

\item Análisis y discusión de resultados
\begin{itemize}
    \item Analizar las ventajas y limitaciones de la metodología propuesta.
    \item Discutir posibles mejoras y aplicaciones del sistema en entornos comerciales o de monitoreo en tiempo real.
\end{itemize}
\end{itemize}
\clearpage

\section{Cronograma}
% \vspace{-20pt}
\begin{table}[H]
	\centering
    \renewcommand{\arraystretch}{2}
	\resizebox{15cm}{!}{
		\begin{tabular}{|p{10cm}||c|c|c|c|c|c|c|c|c|c|c|c|c|c|c|c|c|}
			\hline
			\multicolumn{18}{|c|}{Cronograma de Actividades 2025}\\
			\hline
			Mes & \multicolumn{4}{|c|}{Enero} & \multicolumn{4}{|c|}{Febrero} & \multicolumn{4}{|c|}{Marzo} & \multicolumn{5}{|c|}{Abril}\\
			\hline 
			Actividad/Semana & 2 & 3 & 4 & 5 & 6 & 7 & 8 & 9 & 10 & 11 & 12 & 13  & 14 & 15 & 16 & 17 & 18 \\ 
			\hline
			\hline
			Inicio & X &  &  &  &  &  &  &  &  &  &  &  &  &   &  &  &  \\
			\hline
			Antecedentes, Justificación & X & X & X & X & X & X &  &  &  &  &  &  &  &  &   &  &  \\
			\hline
            Objetivos e Hipótesis &  &  &  &  &  &  & X & X & X &  &  &  &  &  &   &  &  \\
            \hline
            Metodología, Organización de la tesis, Referencias &  &  &  &   &  &  &  &  &  & X & X & X &  &  &   &  &  \\
            \hline            
            Desarrollo teórico del sistema &  &  &  &  &  &  &  &  &  &  &  & & X & X &   &  & X \\
			\hline
             Mes & \multicolumn{4}{|c|}{Mayo} & \multicolumn{4}{|c|}{Junio} & \multicolumn{5}{|c|}{Julio} & \multicolumn{4}{|c|}{Agosto}\\
			\hline 
			Actividad/Semana  & 19 & 20 & 21 & 22 & 23 & 24 & 25 & 26 & 27 & 28 & 29 & 30 & 31 & 32 & 33 & 34 & 35 \\ 
			\hline
		Desarrollo teórico del sistema & X & X &  &  &  &  &  &  &  &  &  & &  &  &   &  &  \\
			\hline
        Simulación del sistema &  &  & X & X & X & X &  &  &  &  &  &  &  &   &  &  & \\
			\hline	
            Implementación de la plataforma de simulación &  &  &  &  &  &  & X & X & X & X & X &  &  &   &  &  & \\
			\hline
			Presentación Anteproyecto &  &  & &  &  &  &  &  &  &  &  & X & X &   &  &  & \\
			\hline
			Diseño del estimador de datos (Filtros) &  &  &  &  &  &  &  &  & &  &  &  &  & X & X & X & X \\
			\hline
                        \hline
            Mes & \multicolumn{5}{|c|}{Septiembre} & \multicolumn{5}{|c|}{Octubre} & \multicolumn{4}{|c|}{Noviembre} & \multicolumn{3}{|c|}{Diciembre}\\
			\hline 
			Actividad/Semana  & 36 & 37 & 38 & 39 & 40 & 40 & 41 & 42 & 43 & 44 & 45 & 46 & 47 & 48 & 49 & 50 & 51 \\ 
			\hline
			Diseño del estimador de datos (Filtros) & X & X & X & X & X &  &  &  & &  &  &  &  & & & & \\
			\hline
            Validación de los resultados (Simulación/Experimental) &  &  &  &  &  & X & X & X & X &  &  &  &  &  &  &  & \\
			\hline
            Redacción de tesis &  &  &  &  &  &  &  & X & X & X & X & X & X & X & X &  &  \\
			\hline
			Examen Cerrado &  &  &  &  &  &  &  &  &  &  &  &  &  &  &  & X &  \\
			\hline
            Mes & \multicolumn{4}{|c|}{Enero} & \multicolumn{4}{|c|}{Febrero} & \multicolumn{4}{|c|}{Marzo} & \multicolumn{5}{|c|}{Abril} \\
            \hline
            Actividad/Semana & 2 & 3 & 4 & 5 & 6 & 7 & 8 & 9 & 10 & 11 & 12 & 13 & 14 & 15 & 16 & 17 & 18 \\
            \hline \hline
            Examen abierto &  &  & X & & & & & & & & & & & & & & \\
            \hline
			 \hline
		\end{tabular}
	}
	\label{actividades}
\end{table}

\begin{thebibliography}{99}
\bibitem{1}
R. Kumaresan y D. W. Tufts, «Estimating the parameters of exponentially damped sinusoids and pole-zero modeling in noise», \emph{IEEE Transactions on Acoustics, Speech, and Signal Processing}, vol. 30, no. 6, pp. 833-840, dic. 1982, doi: 10.1109/TASSP.1982.1163974.

\bibitem{2}
P. Kundur, \emph{Power System Stability and Control}, 1st ed. New York: McGraw-Hill, 1994.

\bibitem{3}
C. A. Ordóñez y M. A. Ríos, «Electromechanical Modes Identification Based on Sliding-window Data from a Wide-area Monitoring System», \emph{Electric Power Components and Systems}, vol. 41, no. 13, pp. 1264-1279, ago. 2013, doi: 10.1080/15325008.2013.816982.

\bibitem{4}
A. Zamora, V. M. Venkatasubramanian, J. A. de la O Serna, J. M. Ramirez, y M. R. A. Paternina, «Multi-dimensional ringdown modal analysis by filtering», \emph{Electric Power Systems Research}, vol. 143, pp. 748-759, feb. 2017, doi: 10.1016/j.epsr.2016.10.016.

\bibitem{5}
E. S. Bañuelos-Cabral, J. J. Nuño-Ayón, B. Gustavsen, J. A. Gutiérrez-Robles, V. A. Galván-Sánchez, J. Sotelo-Castañón, y J. L. García-Sánchez, «Spectral fitting approach to estimate electromechanical oscillation modes and mode shapes by using vector fitting», \emph{Electric Power Systems Research}, vol. 176, p. 105958, nov. 2019, doi: 10.1016/j.epsr.2019.105958.

\bibitem{6}
A. Almunif, L. Fan, y Z. Miao, «A tutorial on data-driven eigenvalue identification: Prony analysis, matrix pencil, and eigensystem realization algorithm», \emph{International Transactions on Electrical Energy Systems}, vol. 30, no. 1, p. e12283, ene. 2020, doi: 10.1002/2050-7038.12283.

\bibitem{7}
X. Li, T. Jiang, H. Yuan, H. Cui, F. Li, G. Li, y H. Jia, «An eigensystem realization algorithm based data-driven approach for extracting electromechanical oscillation dynamic patterns from synchrophasor measurements in bulk power grids», \emph{International Journal of Electrical Power \& Energy Systems}, vol. 116, p. 105549, mar. 2020, doi: 10.1016/j.ijepes.2019.105549.

\bibitem{8}
M. A. Moghdam, R. Noroozian, y S. Jalilzadeh, «Improving the accuracy and the estimation time of inter-area modes in power system based on Tufts-Kumaresan algorithm», \emph{International Journal of Modelling and Simulation}, pp. 1-13, ene. 2023, doi: 10.1080/02286203.2022.2161443.

\bibitem{9}
Y. Chen, F. Wu, L. Shi, Y. Li, P. Qi, y X. Guo, «Identification of Sub-Synchronous Oscillation Mode Based on HO-VMD and SVD-Regularized TLS-Prony Methods», \emph{Energies}, vol. 17, no. 1, p. 187, ene. 2024, doi: 10.3390/en17010187.

\bibitem{10}
J. Ramos, R. Jimenez, E. Barocio «Visualization of Interarea Oscillations Using an Extended Subspace Identification Technique», \emph{North American Power Symposium (NAPS)}, , sep. 2013, doi: 10.1109/NAPS.2013.6666931.

\end{thebibliography}







